% §3 Problem setting.  Point: forming S at the root costs Theta(m^2) pairwise estimates,
% so we observe a Bernoulli(p) subset instead.  Delivers the observation model, the second
% (distance) operator, and the recovery criterion eq:entrywise that every bound is measured
% against.  This file keeps \label{sec:problem}; §2 is \label{sec:model}.
\section{Problem setting}\label{sec:problem}

\subsection{Why a sampled similarity}\label{sec:problem-cost}

\Cref{thm:stdr-split} makes divide-and-conquer over the similarity matrix a viable
strategy for large trees: partition the leaves by the sign of the Fiedler vector,
solve each side independently by any reconstruction method, and merge the two
subtrees. The recursion is attractive because the expensive part of tree
reconstruction is superlinear in the number of leaves, so cutting the problem in
half repeatedly reduces the total work sharply, while \Cref{thm:stdr-split}
guarantees that each cut is a genuine split and the merge is therefore well-posed
\cite{aizenbud2023spectral}.

The top-level call is the one that governs both the error and the cost. It is the
only call that sees all $m$ leaves, and no parent call corrects it: a wrong split at
the root cannot be repaired further down. It is also where the similarity matrix is
largest.

That is the bottleneck. An entry of $S$ is not a table lookup but a statistic of a
pair of sequences: the log-det similarity of \Cref{def:similarity} is estimated
from the joint character counts of leaves $i$ and $j$ over a length-$\ell$
alignment. Forming the full matrix at the top level therefore costs
\begin{equation}\label{eq:cost}
   \binom{m}{2} = \Theta(m^2) \ \text{pairwise estimates},
   \qquad \Theta(m^2\ell) \ \text{character comparisons},
\end{equation}
which for $m \sim 10^5$ leaves is on the order of $10^{10}$ pairwise estimates
before the recursion has taken a single step. The quadratic term is exactly the cost
that divide-and-conquer was introduced to avoid, and it is paid in full at the root.
This raises the question the paper answers: can the top-down split be computed at
almost linear cost, with a guarantee that it is still the same split the full matrix
would have returned?

\subsection{The observation model}\label{sec:problem-sampling}

Our answer is to compute only a random subset of the pairs. Following the
matrix-completion paradigm \cite{candes2009power, candes2008exact}, we observe
entries through a random mask and reweight: let $\Omega\in\{0,1\}^{m\times m}$ be a
symmetric mask whose upper-triangular entries are independent
$\mathrm{Bernoulli}(p)$, $p\in(0,1]$, and form the empirical similarity by
inverse-probability weighting \cite{candes2009matrix},
\begin{equation}\label{eq:sampling_estimator}
   \hat S_{ij}=\begin{cases}\tfrac1p\,\Omega_{ij}\,S_{ij}, & i\ne j,\\ 1, & i=j.\end{cases}
\end{equation}
Only the $\Theta(p m^2)$ pairs in the mask are ever estimated, so $p$ is the
fraction of the cost in \eqref{eq:cost} that is actually paid, and a rate
$p=\Theta(\log m/m)$ would bring the top-level call down to $\Theta(m\log m)$
pairwise estimates.

This estimator is unbiased, $\E[\hat S]=S$, so $\EL=\hat L-L$ is mean zero, where
$\hat L=\hat D-\hat S$. The resulting error is missingness rather than additive
noise: most pairs are left at zero and the rest rescaled by $1/p$, so it is heavy,
structured, and its per-node variance grows as $p\to0$ (\Cref{prop:C-var}).

The scope of the analysis is the population matrix. The empirical log-det estimator
of an individual entry concentrates around $S_{ij}$ at rate
$\mathcal O(1/\sqrt\ell)$ \cite{aizenbud2023spectral}, and the structural
conditions of \Cref{sec:model} are stated as properties of the population $S$,
whose observable proxy is that estimator. 

\subsection{A second operator on the same leaves}\label{sec:problem-distance}

The Laplacian is not the only operator on these leaves whose spectrum carries the
split. A second construction replaces $S$ by the pairwise \emph{distance} matrix
$\Dmat$ --- distinct from the degree matrix $D$ of $L=D-S$ --- built entrywise from
$S$ by $\Dmat_{ij}=-\log S_{ij}$, and centers it, $\Bop:=H\Dmat H$ with
$H=I-\tfrac1m\1\1^\top$. There the split sits in the sign pattern of $\uone$, the
eigenvector of $\Bop$ at the eigenvalue of largest absolute value, rather than in
the Fiedler vector of $L$.

The two operators see the same tree. Under \Cref{asm:gtr,asm:clock} the closed form
\eqref{eq:similarity_gtr} gives, with $d_{ij}$ the path length of
\eqref{eq:mrca_height},
\begin{equation}\label{eq:distance_from_similarity}
   \Dmat_{ij} = -\log S_{ij} = |\mathrm{tr}(Q)|\;d_{ij},
   \qquad \Dmat_{ii}=0,
\end{equation}
so $\Dmat$ is the additive path metric of $\mathcal{T}$ up to the positive scale
$|\mathrm{tr}(Q)|$, and a spectral statement about $\Bop$ is a statement about tree
distances. Since \eqref{eq:distance_from_similarity} is entrywise, the same
$\mathrm{Bernoulli}(p)$ mask $\Omega$ that samples $S$ samples $\Dmat$ unchanged:
both operators are observed under the one model \eqref{eq:sampling_estimator}, and
\Cref{sec:main} states a sampling rate for each. \Cref{app:dist} works out the
population spectrum of $\Bop$ and proves the distance-route rate.


\subsection{The recovery problem}\label{sec:problem-statement}

\medskip
\noindent\textbf{Problem.} Given the sparsely observed estimate $\hat S$, recover the
Fiedler split $\Pisp{S}$ of the full matrix from the sign pattern of the Fiedler vector
$\vvh$ of $\hat L$. We ask for the minimal sampling rate $p$ at which
$\Pisp{\hat S}=\Pisp{S}$ \hp. Since $\E[\hat S]=S$, the question reduces to an
eigenvector-perturbation problem over $\EL=\hat L-L$, which the remainder of the
paper analyzes.

\Cref{thm:stdr-split} already makes $\Pisp{S}$ a clan pair of $\mathcal{T}$, so a rate
at which $\Pisp{\hat S}=\Pisp{S}$ is a rate at which the sub-sampled operator returns a
genuine split of the tree. 


\medskip
Recovery is, at its core, an \emph{entry-wise} event. Two vectors induce the same
sign partition exactly when they agree in sign at every node, so the top-down
partition succeeds precisely when the empirical Fiedler vector agrees in sign with
its population counterpart coordinate by coordinate, that is, when
\begin{equation}\label{eq:entrywise}
   \norm{\vvh-\vv}_\infty<\min_i|\vv_i|.
\end{equation}
Condition \eqref{eq:entrywise} is the formal recovery event that every bound in
the paper is measured against.
